%% kapa: exam=1 round=32 year=2021 period=2 form=A subject=appraisal_law subject_ko=감정평가관계법규 seq=32 number=32
%% source: Q-Net 감정평가사 기출문제 (cst00309) artlSeq=5208397 file=extracted/1008/5208397/A형/2021년도 제32회 감정평가사 1차 2교시 A형.pdf
%% conversion: kordoc 4.24.51 -> tools/gen1.py (auto); review=proofread-sonnet; flags=none
\begin{question}{32}
부동산등기법령상 등기절차에 관한 설명으로 옳은 것은?
\begin{choices}
  \item 법인의 합병에 따른 등기는 등기권리자와 등기의무자가 공동으로 신청하여야 한다.
  \item 등기의무자는 공유물을 분할하는 판결에 의한 등기를 단독으로 신청할 수 없다.
  \item 토지의 분할이 있는 경우에는 그 토지 소유권의 등기명의인은 그 사실이 있는 때부터 1개월 이내에 그 등기를 신청하여야 한다.
  \item 등기관이 직권에 의한 표시변경등기를 하였을 때에는 소유권의 등기명의인은 지체 없이 그 사실을 지적소관청에게 알려야 한다.
  \item 토지가 멸실된 경우에는 그 토지 소유권의 등기명의인은 그 사실이 있는 때부터 14일 이내에 그 등기를 신청하여야 한다.
\end{choices}
\end{question}
\answer{3}
